
\documentclass[11pt,a4paper,oneside]{article}
\usepackage[top=0.85in,bottom=0.85in,left=0.95in,right=0.95in,headheight=14pt]{geometry}
\usepackage[T1]{fontenc}
\usepackage{lmodern}
\usepackage{microtype}
\usepackage{amsmath,amssymb,amsthm,mathtools,bm}
\usepackage{booktabs,array,tabularx,longtable}
\usepackage{xcolor}
\usepackage{enumitem}
\usepackage{fancyhdr}
\usepackage{titlesec}
\usepackage[colorlinks=true,linkcolor=black,urlcolor=blue,citecolor=black]{hyperref}

\setlist{leftmargin=1.4em,itemsep=0.25em,topsep=0.25em}
\titleformat{\section}{\Large\bfseries}{\S\thesection}{0.8em}{}
\titleformat{\subsection}{\large\bfseries}{\thesubsection}{0.8em}{}
\titleformat{\subsubsection}{\normalsize\bfseries}{\thesubsubsection}{0.8em}{}

\newtheorem{definition}{Definition}[section]
\newtheorem{axiom}{Axiom}[section]
\newtheorem{proposition}{Proposition}[section]
\newtheorem{remark}{Remark}[section]
\newtheorem{theorem}{Theorem}[section]

%% ================================================================
%%  MACRO SET --- IKK II v2.0
%%  Harmonised with IKK III v2.0 and IKK IV v1.0
%% ================================================================

%% ---- Identity objects ----
\newcommand{\Ia}{I_a}
\newcommand{\Ib}{I_b}
\newcommand{\Itrue}{\mathfrak{I}_{\mathrm{true}}}
\newcommand{\Ihid}{\mathfrak{I}_{\mathrm{hidden}}}
\newcommand{\Inull}{\mathfrak{I}_{\mathrm{null}}}
\newcommand{\Ihat}{\widehat{\mathfrak{I}}_a}

%% ---- Typed spaces ----
\newcommand{\Ha}{\mathfrak{H}_a}
\newcommand{\Hb}{\mathcal{H}_b}
\newcommand{\Hahat}{\widehat{\mathfrak{H}}_a}
\newcommand{\Kab}{\mathfrak{K}_{ab}}
\newcommand{\Osp}{\mathcal{O}_b}

%% ---- Projection operator anatomy ----
%  Full projection
\newcommand{\Pab}{\mathcal{P}_{a\to b}}
%  Source encoding operator
\newcommand{\Ca}{\mathcal{C}_a}
%  Access operator
\newcommand{\Aab}{\mathcal{A}_{a\to b}}
%  Readout/resolution operator
\newcommand{\Rb}{\mathcal{R}_b}
%  Complement of access (hidden part)
\newcommand{\Qab}{\mathcal{Q}_{a\to b}}
%  Pseudoinverse
\newcommand{\Ppinv}{P_{a\to b}^{+}}
%  Reconstruction/inference operator
\newcommand{\Jba}{\mathcal{J}_{b\to a}}

%% ---- Encoded identity fields ----
\newcommand{\PhiA}{\Phi_a}              % full encoded field
\newcommand{\PhiAb}{\Phi_{a|b}}         % accessible portion
\newcommand{\PhiHid}{\Phi_{\mathrm{hid}}}  % hidden portion

%% ---- Distinguishability measures ----
\newcommand{\D}{\mathfrak{D}}
\newcommand{\Da}{\mathfrak{D}_a}
\newcommand{\Db}{\mathfrak{D}_b}
\newcommand{\Dhid}{\mathfrak{D}_{\mathrm{hid}}}
\newcommand{\Dacc}{\mathfrak{D}_{\mathrm{acc}}}
\newcommand{\Rhid}{\mathfrak{R}_{a|b}}

%% ---- Efficiency ----
\newcommand{\etaab}{\eta_{ab}}
\newcommand{\etak}{\eta_k(w)}

%% ---- w coordinates ----
\newcommand{\wa}{w_a}
\newcommand{\wb}{w_b}
\newcommand{\Dab}{\Delta_{ab}}

%% ---- Wave state ----
\newcommand{\psiab}{\psi_b}
\newcommand{\Wab}{\mathcal{W}_{a\to b}}

%% ---- Other shorthands ----
\newcommand{\Oh}{\mathcal{O}}
\newcommand{\Loss}{\mathcal{L}_\Pi}
\newcommand{\Rel}[2]{R_{#1#2}}
\newcommand{\MM}{\mathbb{M}}
\newcommand{\RR}{\mathbb{R}}
\newcommand{\CC}{\mathbb{C}}
\newcommand{\TT}{\mathbb{T}}
\newcommand{\wcoord}{w}

%% ================================================================

\pagestyle{fancy}
\fancyhf{}
\fancyhead[L]{\small\textit{Inverse Kaluza-Klein Framework --- IKK Paper Set}}
\fancyhead[R]{\small\textit{IKK II v2.0}}
\fancyfoot[C]{\thepage}
\renewcommand{\headrulewidth}{0.4pt}
\setlength{\parskip}{0.55em}
\setlength{\parindent}{0pt}

\begin{document}

\begin{titlepage}
\centering
\vspace*{1.1cm}
{\large\scshape VSEPR-SIM Theoretical Division / IKK Paper Set\par}
\vspace{1.1cm}
{\Huge\bfseries IKK II\par}
\vspace{0.35cm}
{\Large\bfseries Projection, Observable Form, and Residual Identity\par}
\vspace{0.2cm}
{\large Version 2.0\par}
\vspace{1.1cm}
{\normalsize
This document upgrades the projection foundation of the IKK framework.
The projection operator is given explicit anatomy:
$\Pab = \Rb\,\Aab\,\Ca$.
The source encoding operator $\Ca$ prepares primitive identity for projection.
The access operator $\Aab(\wa, \wb, \Dab)$ determines what fraction of
encoded identity is accessible from observation domain $b$,
giving the fifth coordinate $\wcoord$ an explicit operational role.
The readout operator $\Rb$ maps accessible identity to observable form.
Distinguishability is defined as a norm ratio,
the Born rule is derived as normalized observed distinguishability,
and hidden content is identified as the projection-kernel content of identity.
\par}
\vfill
{\large Companion: IKK I (foundations), IKK III (matrix/wave bridge), IKK IV (scale intervals)\par}
\end{titlepage}

\tableofcontents
\clearpage

%% ================================================================
\section{Scope and Upgrade Summary}
%% ================================================================

IKK II is the projection-foundation document of the IKK paper set. Version~2 retains the conceptual spine of v1 --- existence gate, inverse inference, typed spaces, hidden residual --- and adds the formal construction that was previously absent:

\[
\boxed{\Pab = \Rb\,\Aab\,\Ca.}
\]

This gives the projection operator \emph{anatomy}. The operator is no longer a bare arrow; it has three identifiable components with distinct mathematical roles. The central consequences:

\begin{enumerate}
\item Distinguishability is defined as a norm ratio, not a vague ordering.
\item The coordinate $\wcoord$ acquires an explicit operational role inside $\Aab$.
\item The Born rule follows as normalized observed distinguishability.
\item Hidden content is exactly the projection-kernel content of the identity state.
\item A reconstruction operator $\Jba = \Ppinv$ is defined, with its null-space gap made precise.
\end{enumerate}

\textbf{What moved:} dark matter kernels, proper-time dynamics, relation-particle chains, and code prototypes remain in later documents.

%% ================================================================
\section{Central Thesis}
%% ================================================================

\[
\boxed{O_b = \Pab[\Ia] = \Rb\,\Aab\,\Ca[\Ia],\qquad O_b \neq \Ia.}
\]

The observable $O_b$ is not the identity object $\Ia$. It is the result of encoding $\Ia$ into a projectable field, filtering that field through an access operator parameterized by $\wcoord$, and reading out the accessible portion in the target domain $b$.

Every equation in IKK~II follows from this factorization.

%% ================================================================
\section{Inverse Direction of Inference}
%% ================================================================

Standard Kaluza-Klein infers observable force behavior from an assumed extra dimension:
\[
\text{KK:}\quad D_{\mathrm{extra}} \longrightarrow F_{\mathrm{obs}}.
\]

IKK reverses this:
\[
\boxed{\text{IKK:}\quad F_{\mathrm{obs}}(S) \longrightarrow D_{\mathrm{perm}}(S).}
\]

$D_{\mathrm{perm}}$ is a permission structure: a hidden-access condition required for the observed force-scale behavior to be representable. IKK begins with observed residuals or representation loss and asks what access structure must exist for such observations to be recoverable.

%% ================================================================
\section{Existence Gate}
%% ================================================================

\begin{axiom}[Existence gate]
Every identity-bearing object satisfies
$E \in \{0,1\}.$
$E = 1$ admits the object into the identity accounting.
It does not yet specify visibility, measurability, or projection accessibility.
\end{axiom}

\[
E(\Ia) = 1 \;\not\Rightarrow\; \Ia = O_b.
\]

Existence is separated from observability. An object may exist and produce no observable output under a given projection basis.

%% ================================================================
\section{Typed Identity Spaces}
%% ================================================================

\begin{definition}[Source identity space]
$\Ha$ is the source identity space carrying the full identity content of a system before projection. Objects in $\Ha$ are not in general directly readable as Hilbert-space states.
\end{definition}

\begin{definition}[Observable space]
$\Osp$ is the target observable space. Objects in $\Osp$ are readable under the projection basis.
\end{definition}

The projection operator maps between them:
\[
\Pab : \Ha \longrightarrow \Osp.
\]

$\Ha$ is not assumed to be a Hilbert space. The structure of $\Ha$ constrains what $\Ca$ can do; the structure of $\Osp$ constrains what $\Rb$ can produce.

%% ================================================================
\section{Projection Operator: Full Anatomy}
%% ================================================================

The projection operator has three components.

\[
\boxed{\Pab = \Rb\,\Aab\,\Ca}
\]

Each component has a distinct role.

\subsection{Component 1: Source Encoding Operator $\Ca$}

$\Ca$ prepares the primitive identity for projection. It maps $\Ia$ into an encoded identity field $\PhiA$ that lives in an intermediate space of dimension $N$:

\begin{definition}[Source encoding]
\[
\Ca : \Ia \mapsto \PhiA,\qquad \PhiA = \Ca[\Ia].
\]
$\Ca \in \mathbb{C}^{N \times n}$ in the finite matrix model, where $\Ia \in \mathbb{C}^n$ and $N \geq n$.
\end{definition}

This matters because the primitive identity $\Ia$ may not directly live in the same mathematical space as the observable representation. The encoding lifts $\Ia$ into an intermediate field space where projection can act.

\begin{remark}
$\Ia \neq \psiab$ in general. But $\Ca[\Ia]$ may produce something wave-like after encoding. The wavefunction is not the primitive identity; it is a post-encoding, post-projection readout.
\end{remark}

\subsection{Component 2: Access Operator $\Aab$}

$\Aab$ determines what fraction of the encoded field $\PhiA$ is accessible from observation domain $b$:

\begin{definition}[Access operator]
\[
\Aab : \PhiA \mapsto \PhiAb,\qquad \PhiAb = \Aab[\PhiA].
\]
$\PhiAb$ is the portion of source identity accessible from domain $b$.
\end{definition}

The access operator depends on the projection coordinates of the source and observer and on the domain mismatch:

\[
\boxed{\Aab = \mathcal{A}(\wa,\, \wb,\, \Dab)}
\]

where:
\begin{itemize}
\item $\wa$ = source projection coordinate;
\item $\wb$ = observer/probe projection coordinate;
\item $\Dab$ = scale/domain mismatch between $a$ and $b$.
\end{itemize}

This makes $\wcoord$ \emph{operational}. It is not a spatial axis. It is not a length scale. It controls which channels of encoded identity survive projection into a given observation domain.

\subsection{Component 3: Readout Operator $\Rb$}

$\Rb$ maps accessible identity into the observable representation:

\begin{definition}[Readout operator]
\[
\Rb : \PhiAb \mapsto O_b,\qquad O_b = \Rb[\PhiAb].
\]
$\Rb \in \mathbb{C}^{m \times N}$ in the finite matrix model.
\end{definition}

\subsection{The Factored Equation}

Combining:
\[
\boxed{O_b = \Pab[\Ia] = \Rb\,\Aab\,\Ca[\Ia].}
\]

This is the central equation of IKK~II v2. The projection has anatomy: encoding, access filtering, readout.

%% ================================================================
\section{Matrix Version}
%% ================================================================

For a finite model with $\Ia \in \mathbb{C}^n$, $O_b \in \mathbb{C}^m$, and intermediate dimension $N$:

\[
\Ca \in \mathbb{C}^{N \times n},\qquad
\Aab \in \mathbb{C}^{N \times N},\qquad
\Rb \in \mathbb{C}^{m \times N}.
\]

\[
\boxed{P_{a\to b} = R_b\, A_{a\to b}\, C_a \;\in\; \mathbb{C}^{m \times n}.}
\]

The projected observable amplitude:
\[
\psiab = P_{a\to b}\,\Ia = R_b\, A_{a\to b}\, C_a\,\Ia.
\]

This is now computable. Given explicit matrices $C_a$, $A_{a\to b}$, $R_b$, and an identity vector $\Ia$, the observable $O_b$ is a matrix product. The framework is no longer purely formal.

%% ================================================================
\section{Distinguishability as a Norm}
%% ================================================================

\begin{definition}[Source distinguishability]
\[
\Da(\Ia) = \|\Ca\,\Ia\|^2.
\]
\end{definition}

\begin{definition}[Observed distinguishability]
\[
\Db(\Ia) = \|\Pab\,\Ia\|^2 = \|\Rb\,\Aab\,\Ca\,\Ia\|^2.
\]
\end{definition}

\begin{definition}[Projection efficiency]
\[
\boxed{\etaab = \frac{\|\Rb\,\Aab\,\Ca\,\Ia\|^2}{\|\Ca\,\Ia\|^2} = \frac{\Db(\Ia)}{\Da(\Ia)}.}
\]
\end{definition}

$\etaab \in [0,1]$. The efficiency is now not a declaration but a ratio of norms. When $\Aab = I$ and $\Rb\,\Ca$ is unitary, $\etaab = 1$.

This is the first version of $\etaab = \Db / \Da$ that actually means something. It says: $\etaab$ is the fraction of source identity distinguishability retained after encoding, filtering, and readout.

%% ================================================================
\section{Projection Decomposition}
%% ================================================================

Define the complement of the access operator:
\[
\Qab = I - \Aab.
\]

Then the hidden portion of the encoded field is:
\[
\PhiHid = \Qab\,\Ca\,\Ia.
\]

\begin{proposition}[Orthogonal decomposition]
If $\Aab$ is an orthogonal projector ($\Aab^2 = \Aab$, $\Aab^\dagger = \Aab$), then:
\[
\|\PhiA\|^2 = \|\Aab\,\PhiA\|^2 + \|(I - \Aab)\,\PhiA\|^2.
\]
\end{proposition}

This gives a bookkeeping law:

\[
\boxed{\Da = \Dacc + \Dhid}
\]

where $\Dacc = \|\Aab\,\Ca\,\Ia\|^2$ and $\Dhid = \|\Qab\,\Ca\,\Ia\|^2$.

Equivalently:
\[
\boxed{\Da = \Db + \Rhid}
\]

where $\Rhid = \|\Qab\,\Ca\,\Ia\|^2$ is the hidden residual distinguishability. Hidden content is not a free parameter. It is constrained by $\Rhid = \Da - \Db \geq 0$.

This prevents the hidden residual from becoming a theoretical escape clause. No free excuses.

%% ================================================================
\section{The Access Operator and the Role of $w$}
%% ================================================================

In the diagonal basis of the encoded identity channels, define per-channel access factors:

\begin{definition}[Diagonal access operator]
\[
\Aab(w) = \mathrm{diag}\!\left(\sqrt{\eta_1(w)},\,\sqrt{\eta_2(w)},\,\ldots,\sqrt{\eta_N(w)}\right)
\]
where $0 \leq \eta_k(w) \leq 1$ for each channel $k$.
\end{definition}

The encoded field component at channel $k$ survives with amplitude:
\[
\Phi_{a|b,\,k} = \sqrt{\etak}\;\Phi_{a,k}.
\]

$\wcoord$ controls how much of each channel survives projection. Two candidate forms:
\[
\eta_k(w) = \exp\!\left[-\lambda_k(\wa - \wb)^2\right]
\]
\[
\eta_k(w) = \frac{1}{1 + \lambda_k(\wa - \wb)^2}
\]

where $\lambda_k$ controls the sensitivity of channel $k$ to projection mismatch. When $\wa = \wb$, every channel survives. When $|\wa - \wb| \gg \lambda_k^{-1/2}$, channel $k$ is suppressed.

\begin{definition}[Operational role of $w$]
\[
\boxed{w \text{ controls channel-wise access to source identity.}}
\]
$w$ is the coordinate governing projective accessibility between the source identity domain and the observation domain.
\end{definition}

This is not "length scale." It is not "another spatial dimension." It is the access parameter embedded in $\Aab$. This definition is parallel to, and made explicit by, the interval-level $\chi_w$ construction in IKK~IV.

%% ================================================================
\section{Wavefunction as Projection Output}
%% ================================================================

The wavefunction is not primitive. It is the readout form of accessible identity:

\[
\boxed{\psiab = \Rb\,\Aab(w)\,\Ca\,\Ia.}
\]

Written through the wave-state recovery map $\Wab = \Rb \circ \Aab \circ \Ca : \Ha \to \Hb$:
\[
\psiab = \Wab(\Ia) \in \Hb.
\]

The core interpretive statement:
\begin{quote}
The wavefunction is the projected observable amplitude produced when primitive identity $\Ia$ is encoded by $\Ca$, filtered for accessibility by $\Aab(w)$, and read out by $\Rb$. It is not the identity object. It is the portion of encoded identity that survives projection into the observation domain.
\end{quote}

This is the wave-particle bridge in one sentence. The wavefunction is real, useful, and predictive precisely because $\Hb$ preserves the quantities that quantum mechanics needs. It is not a falsehood. It is a projection.

%% ================================================================
\section{Born Rule as Normalized Distinguishability}
%% ================================================================

In the finite model, let $i$ index the output channels of $\Rb$. Measurement probability in channel $i$ is:

\[
\boxed{p_i = \frac{\left|(R_b\,A_{a\to b}\,C_a\,\Ia)_i\right|^2}{\displaystyle\sum_j \left|(R_b\,A_{a\to b}\,C_a\,\Ia)_j\right|^2}.}
\]

The Born rule is not "wavefunction magic." It is normalized observed distinguishability. The probability of observing outcome $i$ is the fraction of total observed distinguishability carried by channel $i$.

\begin{proposition}[Born rule recovery]
When $\Ca = I$, $\Rb = I$, and $\Aab = \sqrt{\eta}\,I$ (uniform access), the above reduces to:
\[
p_i = \frac{|(\Ia)_i|^2}{\sum_j |(\Ia)_j|^2},
\]
which is the standard Born rule under normalization $\|\Ia\| = 1$.
\end{proposition}

%% ================================================================
\section{Minimal Example: Two-Channel Identity}
%% ================================================================

Let $\Ia = [\alpha,\,\beta]^T$, $\Ca = I$, $\Rb = I$, and:
\[
A_{a\to b}(w) = \begin{bmatrix} \sqrt{\eta_1} & 0 \\ 0 & \sqrt{\eta_2} \end{bmatrix}.
\]

Then:
\[
\psiab = \begin{bmatrix} \sqrt{\eta_1}\,\alpha \\ \sqrt{\eta_2}\,\beta \end{bmatrix}.
\]

Observed distinguishability by channel:
\[
\mathfrak{D}_{b,1} = \eta_1|\alpha|^2,\qquad \mathfrak{D}_{b,2} = \eta_2|\beta|^2.
\]

Measurement probabilities:
\[
\boxed{p_1 = \frac{\eta_1|\alpha|^2}{\eta_1|\alpha|^2 + \eta_2|\beta|^2},\qquad p_2 = \frac{\eta_2|\beta|^2}{\eta_1|\alpha|^2 + \eta_2|\beta|^2}.}
\]

When $\eta_1 = \eta_2$, the $\eta$ factors cancel and the standard Born probabilities $|\alpha|^2$ and $|\beta|^2$ are recovered (assuming normalization). When $\eta_1 \neq \eta_2$, the probabilities are weighted by differential accessibility. This is the first deformation equation: a prediction of modified measurement statistics when projection accessibility is channel-dependent.

%% ================================================================
\section{The Three Projection Regimes}
%% ================================================================

\paragraph{Regime I: Perfect projection.}
$\Aab = I$. Full accessibility.
\[
\psiab = \Ca\,\Ia,\qquad \etaab = 1,\qquad \Rhid = 0.
\]

\paragraph{Regime II: Uniform lossy projection.}
$\Aab = \sqrt{\eta}\,I$ for scalar $\eta \in (0,1)$.
\[
\psiab = \sqrt{\eta}\,\Ca\,\Ia.
\]
Relative probabilities recover the standard Born rule after normalization. Total detectability is reduced. All channels are suppressed equally. This corresponds to weak coupling or symmetric scale mismatch.

\paragraph{Regime III: Channel-dependent projection.}
$\Aab = \mathrm{diag}(\sqrt{\eta_1},\ldots,\sqrt{\eta_N})$ with $\eta_i \neq \eta_j$.
\[
p_i = \frac{\eta_i\,|\psi_i|^2}{\sum_j \eta_j\,|\psi_j|^2}.
\]

Regime III predicts deviations from ordinary quantum measurement statistics. This is where EVM begins: testing whether physical systems exhibit channel-dependent access and whether the deformation follows the predicted $\eta_k(w)$ form.

%% ================================================================
\section{Hidden Residual and the Complement Operator}
%% ================================================================

\begin{definition}[Hidden residual field]
\[
\PhiHid = \Qab\,\Ca\,\Ia = (I - \Aab)\,\Ca\,\Ia.
\]
\end{definition}

The hidden residual $\PhiHid$ is not destroyed content. It is the complement of $\PhiAb$ under $\Aab$. Its norm is:
\[
\Dhid = \|\Qab\,\Ca\,\Ia\|^2 = \Da - \Dacc.
\]

By the bookkeeping law (\S\ref{sec:decomp}):
\[
0 \leq \Rhid \leq \Da.
\]

The hidden residual is finite and bounded. It cannot exceed the total source distinguishability. This is the formal version of "hidden residual is not a magic trash bin."

\label{sec:decomp}

%% ================================================================
\section{Hidden Variables as Projection-Kernel Content}
%% ================================================================

\begin{definition}[Projection kernel]
\[
\ker(\Pab) = \{\Ia \in \Ha : \Pab\,\Ia = 0\}.
\]
\end{definition}

\begin{definition}[Hidden content]
Identity content in $\ker(\Pab)$ produces no observable output under projection $\Pab$.
\[
\Ihid \in \ker(\Pab).
\]
\end{definition}

\[
\boxed{\text{Hidden} = \text{identity content in the projection kernel of the current observation map.}}
\]

This is a precise statement about hidden variables. It is not "there might be hidden variables somewhere." It is: hidden identity content is specifically and only the content lying in $\ker(\Pab)$.

The kernel also resolves the temporal distinction introduced in IKK~II v1:

\begin{itemize}
\item $H_t$: hidden but time-dependent --- content in $\ker(P^{(x,y,z)})$ but not in $\ker(P^{(x,y,z,t)})$.
\item $H_{\neg t}$: hidden and time-independent --- content in $\ker(P^{(x,y,z,t)})$.
\end{itemize}

In symbolic form:
\[
\ker(P_{\mathrm{space}}) \neq \ker(P_{\mathrm{spacetime}}).
\]

A system may be temporally observable but spatially hidden, or hidden from both. The projection map determines the answer. This lets WIC handle both interpretations with the same formalism.

%% ================================================================
\section{Reconstruction Operator}
%% ================================================================

Given $O_b = \Pab\,\Ia$, define the reconstruction operator:
\[
\Jba : \Osp \longrightarrow \Hahat,\qquad \Ihat = \Jba[O_b] = \Ppinv\,O_b.
\]

where $\Ppinv$ denotes the Moore-Penrose pseudoinverse of $P_{a\to b}$.

The pseudoinverse does not recover $\Ia$ exactly. It recovers only the component of $\Ia$ in the orthogonal complement of $\ker(\Pab)$:
\[
\Ppinv\,P_{a\to b}\,\Ia = \Pi_{\mathrm{rec}}\,\Ia.
\]

Consequently:
\[
\boxed{\Ia = \Ihat + \Inull,\qquad \Inull \in \ker(\Pab).}
\]

$\Ihat$ is the recoverable portion. $\Inull$ is the null-space component --- exactly the hidden content. The gap between $\Ia$ and $\Ihat$ is not failure of the reconstruction; it is the mathematically necessary consequence of projection loss.

\begin{remark}
The null-space gap is not a defect of the framework. It is the formal statement that projection is lossy. Any framework claiming perfect recovery from incomplete projection data is wrong. This one does not.
\end{remark}

%% ================================================================
\section{Matrix Mechanics and Wave Mechanics as Parallel Projections}
%% ================================================================

Both matrix mechanics and wave mechanics emerge from the same factored projection with different choices of $\Ca$, $\Aab$, $\Rb$:

\[
\Ia \xrightarrow{\;P_M\;} M \in \Hb \qquad \text{(matrix projection: } C_a^M,\, A^M,\, R_b^M\text{)}
\]
\[
\Ia \xrightarrow{\;P_\psi\;} \psiab \in \Hb \qquad \text{(wave projection: } C_a^\psi,\, A^\psi,\, R_b^\psi\text{)}
\]

The two are equivalent at the level of $\Hb$ predictions when the two projection chains factor through the same intermediate space $\Kab$ with compatible bases. IKK~III provides the full typed bridge $\Wab = \Rb \circ \Aab \circ \Ca$ and the conditions under which $M$ and $\psiab$ produce identical Born-rule outputs.

%% ================================================================
\section{MIR Bridge}
%% ================================================================

IKK~II provides the projection ontology. The MIR paper set provides representation-quality scoring.

The three representation domains enter here only as a bridge:
\[
\mathbb{M}_\RR(S),\qquad \mathbb{M}_\CC(S),\qquad \mathbb{M}_\TT(S).
\]

IKK~II establishes: $O_b = \Rb\,\Aab\,\Ca\,\Ia$ and $\etaab = \Db/\Da$. MIR ranks representations by how well they encode and recover identity under these projections. The argument that efficient storage proves truth is not a scientific argument. Efficiency is evidence, not ontology.

%% ================================================================
\section{Forward References}
%% ================================================================

\paragraph{Dark matter residuals.} Projection residuals at galactic scales: $R_{\mathrm{gal}} = \mathcal{D}_{\mathrm{obs}} - \mathcal{D}_{\mathrm{visible}}$. Forward reference: IKK~V, EVM~I.

\paragraph{Proper-time extension.} Two systems with identical $(x,y,z,t)$ trajectories may differ in recovered identity history if $\wcoord$ evolves differently. Forward reference: IKK~IV \S6.

\paragraph{Relation identity.} $\Rel{A}{B} \neq A + B$; the relation is a distinct element. Forward reference: WIC~III.

%% ================================================================
\section{Overclaim Guardrails}
%% ================================================================

\begin{enumerate}
\item Projection formalism does not prove a hidden ontology. The formal structure requires empirical testing of $\eta_k(w)$.
\item Wavefunction-as-projection must recover standard QM in the limit $\eta_1 = \eta_2 = \cdots = \eta_N$. If it does not, the framework fails.
\item Regime III (channel-dependent access) is a prediction, not an explanation. It must be tested before being used to explain anomalies.
\item Hidden residuals have a bookkeeping law: $\Rhid = \Da - \Db \geq 0$. Claimed residuals must respect it.
\item The null-space gap $\Inull$ is fixed by $\ker(\Pab)$; it is not a free variable.
\end{enumerate}

%% ================================================================
\section{Abstract and Thesis}
%% ================================================================

\textbf{Abstract.} This document develops the projection anatomy of the Inverse Kaluza-Klein framework. The projection operator is factored as $\Pab = \Rb\,\Aab\,\Ca$, giving each component an explicit mathematical role. The source encoding operator $\Ca$ lifts primitive identity into an intermediate field space. The access operator $\Aab(\wa,\wb,\Dab)$ gives the fifth coordinate $\wcoord$ an operational definition: it controls channel-wise accessibility between source and observation domain. The readout operator $\Rb$ maps accessible field content to the observable space. Distinguishability is defined as a norm ratio $\etaab = \|\Pab\,\Ia\|^2/\|\Ca\,\Ia\|^2$. The Born rule follows as normalized observed distinguishability. Channel-uniform access recovers standard quantum probability; channel-dependent access predicts a deformed Born rule testable by the EVM series. Hidden content is defined as projection-kernel content, giving hidden variables a precise and bounded meaning.

\[
\boxed{O_b = \Rb\,\Aab(w)\,\Ca\,\Ia,\qquad \etaab = \frac{\Db}{\Da},\qquad \Ia = \Ihat + \Inull,\quad \Inull \in \ker(\Pab).}
\]

%% ================================================================
\section{Document Structure Map}
%% ================================================================

\begin{longtable}{p{0.10\linewidth}p{0.32\linewidth}p{0.48\linewidth}}
\toprule
Pages & Section & Function \\
\midrule
1--2   & Scope; upgrade summary       & What v2 adds over v1. \\
2--3   & Central thesis; inference     & $O_b = R_b A C_a I_a$; IKK reversal. \\
3--4   & Existence gate; typed spaces  & $E=1$; $\Ha$, $\Osp$ defined. \\
4--7   & Projection anatomy            & $C_a$, $A_{a\to b}$, $R_b$; matrix version. \\
7--9   & Distinguishability as norm    & $\eta_{ab}$; decomposition; bookkeeping law. \\
9--11  & Access operator and $w$       & Diagonal $A$; $\eta_k(w)$ forms; $w$ defined. \\
11--13 & Wavefunction; Born rule       & $\psi_b = R_b A C_a I_a$; normalized $\mathfrak{D}$. \\
13--14 & Two-channel example           & Explicit $2\times 2$ derivation; deformation. \\
14--15 & Three regimes                 & Perfect, uniform lossy, channel-dependent. \\
15--17 & Hidden residual; kernel       & $\mathfrak{Q}$; bookkeeping; $\ker(P)$ = hidden. \\
17--18 & Reconstruction operator       & $J_{b\to a} = P^+$; null-space gap. \\
18--20 & Matrix/wave parallel; MIR    & Two projection chains; representation scoring. \\
20--21 & Forward refs; guardrails      & IKK V, WIC III, IKK IV; five limits. \\
21--22 & Abstract; map                 & Thesis and structure. \\
\bottomrule
\end{longtable}

\end{document}
